% GENERATED FILE. Edit canonical lessons, then run npm run generate:books.
% canonical-source-hash: fnv1a64:06cdc60452c4c2ab
% canonical-lessons: GU-C183-grahak, GU-C183-dukandar, GU-C183-kariyanu, GU-C183-tolvu, GU-C183-kilo

\chapter{At the grocer's counter}
\label{ch:gu-at-the-grocer-s-counter}
\hlchaptermodality{183}

\begin{chapteropening}
\textbf{By the end of this chapter:} \emph{I can name the people and the goods at a grocer's shop.}

Complete the last lesson of chapter 183: I can name the people and the goods at a grocer's shop.
\end{chapteropening}

\section[\texorpdfstring{\gu{ગ્રાહક} (grāhak)}{ગ્રાહક (grāhak)}]{\gu{ગ્રાહક} (grāhak) — customer, buyer}

\label{lesson:GU-C183-grahak}



\begin{quote}
Before the new one: say the Gujarati for to believe, then the Gujarati for to change.
\end{quote}

\subsection*{You'll want to know: \gu{ગ્રાહક}}
\textbf{\gu{ગ્રાહક}} — \emph{grāhak} — \textquotedblleft{}customer, buyer\textquotedblright{}.

\begin{quote}
\textbf{\gu{દુકાનમાં} \gu{એક} \gu{ગ્રાહક} \gu{છે}.} — \emph{dukānmā̃ ek grāhak che} — \textquotedblleft{}There is one customer in the shop.\textquotedblright{}
\end{quote}

A masculine noun.

\begin{cousinweb}
\textbf{\gu{ગ્રાહક}} (\emph{grāhak}) is Sanskrit \emph{grāhaka}, \textquotedblleft{}one who takes, a buyer\textquotedblright{}, taken over whole as a learned loan, a tatsama. It is built on the root \emph{grah-}, \textquotedblleft{}to seize, to take\textquotedblright{}, which CDIAL lists in its causative \emph{grāhayati}, \textquotedblleft{}causes to be taken\textquotedblright{} (CDIAL 4383).
\end{cousinweb}

\begin{grammarlens}[title={What you've built}]
One more word for everyday life at A2.
\end{grammarlens}

\subsection*{Your turn}
\begin{itemize}
\raggedright
  \item \emph{Say it:} \emph{grāhak}
  \item \emph{Say it:} \emph{grāhak}, once more
  \item \emph{Say it:} \emph{dukānmā̃ ek grāhak che} — There is one customer in the shop.
  \item \emph{Recall:} say the Gujarati for to believe, then the Gujarati for to change, then say \emph{grāhak} again
\end{itemize}

\begin{tcolorbox}[breakable,colback=teal!4,colframe=teal!35!black,title={Before you move on}]
What does \gu{ગ્રાહક} mean? (\textquotedblleft{}Customer, buyer\textquotedblright{}.) Say it once more.
\end{tcolorbox}
\section[\texorpdfstring{\gu{દુકાનદાર} (dukāndār)}{દુકાનદાર (dukāndār)}]{\gu{દુકાનદાર} (dukāndār) — shopkeeper}

\label{lesson:GU-C183-dukandar}



\begin{quote}
Before the new one: say the Gujarati for to change, then the Gujarati for customer.
\end{quote}

\subsection*{You'll want to know: \gu{દુકાનદાર}}
\textbf{\gu{દુકાનદાર}} — \emph{dukāndār} — \textquotedblleft{}shopkeeper\textquotedblright{}.

\begin{quote}
\textbf{\gu{આ} \gu{દુકાનદાર} \gu{સારો} \gu{છે}.} — \emph{ā dukāndār sāro che} — \textquotedblleft{}This shopkeeper is a good one.\textquotedblright{}
\end{quote}

Built on \textbf{\gu{દુકાન}}, \emph{dukān}, 'shop'.

\begin{cousinweb}
\textbf{\gu{દુકાનદાર}} (\emph{dukāndār}) is built from \emph{dukān}, \textquotedblleft{}shop\textquotedblright{}, a loan that reached Gujarati through Persian, and the Persian suffix \emph{-dār}, \textquotedblleft{}holder, keeper\textquotedblright{}. Lane lists Arabic \emph{dukkān}, \textquotedblleft{}a shop\textquotedblright{}, under the root \emph{d-k-n}, but notes that some Arabic scholars held it to be a Persian word taken into Arabic, so its first source is not settled (Lane).
\end{cousinweb}

\begin{grammarlens}[title={What you've built}]
One more word for everyday life at A2.
\end{grammarlens}

\subsection*{Your turn}
\begin{itemize}
\raggedright
  \item \emph{Say it:} \emph{dukāndār}
  \item \emph{Say it:} \emph{dukāndār}, once more
  \item \emph{Say it:} \emph{ā dukāndār sāro che} — This shopkeeper is a good one.
  \item \emph{Recall:} say the Gujarati for to change, then the Gujarati for customer, then say \emph{dukāndār} again
\end{itemize}

\begin{tcolorbox}[breakable,colback=teal!4,colframe=teal!35!black,title={Before you move on}]
What does \gu{દુકાનદાર} mean? (\textquotedblleft{}Shopkeeper\textquotedblright{}.) Say it once more.
\end{tcolorbox}
\section[\texorpdfstring{\gu{કરિયાણું} (kariyāṇũ)}{કરિયાણું (kariyāṇũ)}]{\gu{કરિયાણું} (kariyāṇũ) — groceries}

\label{lesson:GU-C183-kariyanu}



\begin{quote}
Before the new one: say the Gujarati for customer, then the Gujarati for shopkeeper.
\end{quote}

\subsection*{You'll want to know: \gu{કરિયાણું}}
\textbf{\gu{કરિયાણું}} — \emph{kariyāṇũ} — \textquotedblleft{}groceries\textquotedblright{}.

\begin{quote}
\textbf{\gu{તે} \gu{બજારથી} \gu{કરિયાણું} \gu{લાવે} \gu{છે}.} — \emph{te bajārthī kariyāṇũ lāve che} — \textquotedblleft{}She brings the groceries from the market.\textquotedblright{}
\end{quote}

A neuter noun.

\begin{cousinweb}
\textbf{\gu{કરિયાણું}} (\emph{kariyāṇũ}) comes from Sanskrit \emph{krayāṇaka}, \textquotedblleft{}goods for sale\textquotedblright{}, built on \emph{krī}, \textquotedblleft{}to buy\textquotedblright{}. CDIAL lists the Gujarati word as \emph{kiriyāṇũ}, \textquotedblleft{}groceries, drugs\textquotedblright{}, among the forms it marks as loan-words (CDIAL 3584).
\end{cousinweb}

\begin{grammarlens}[title={What you've built}]
One more word for everyday life at A2.
\end{grammarlens}

\subsection*{Your turn}
\begin{itemize}
\raggedright
  \item \emph{Say it:} \emph{kariyāṇũ}
  \item \emph{Say it:} \emph{kariyāṇũ}, once more
  \item \emph{Say it:} \emph{te bajārthī kariyāṇũ lāve che} — She brings the groceries from the market.
  \item \emph{Recall:} say the Gujarati for customer, then the Gujarati for shopkeeper, then say \emph{kariyāṇũ} again
\end{itemize}

\begin{tcolorbox}[breakable,colback=teal!4,colframe=teal!35!black,title={Before you move on}]
What does \gu{કરિયાણું} mean? (\textquotedblleft{}Groceries\textquotedblright{}.) Say it once more.
\end{tcolorbox}
\section[\texorpdfstring{\gu{તોલવું} (tolvũ)}{તોલવું (tolvũ)}]{\gu{તોલવું} (tolvũ) — to weigh}

\label{lesson:GU-C183-tolvu}



\begin{quote}
Before the new one: say the Gujarati for shopkeeper, then the Gujarati for groceries.
\end{quote}

\subsection*{You'll want to know: \gu{તોલવું}}
\textbf{\gu{તોલવું}} — \emph{tolvũ} — \textquotedblleft{}to weigh\textquotedblright{}.

\begin{quote}
\textbf{\gu{દુકાનદાર} \gu{ચોખા} \gu{તોલે} \gu{છે}.} — \emph{dukāndār cokhā tole che} — \textquotedblleft{}The shopkeeper weighs the rice.\textquotedblright{}
\end{quote}

\begin{cousinweb}
\textbf{\gu{તોલવું}} (\emph{tolvũ}) is a tadbhava, a word inherited from Sanskrit and reshaped by sound change. It comes from \emph{tōláyati}, \textquotedblleft{}lifts, weighs, considers\textquotedblright{}, and CDIAL lists the Gujarati verb as \emph{toḷvũ}, \textquotedblleft{}to weigh\textquotedblright{} (CDIAL 5979).
\end{cousinweb}

\begin{grammarlens}[title={What you've built}]
One more word for everyday life at A2.
\end{grammarlens}

\subsection*{Your turn}
\begin{itemize}
\raggedright
  \item \emph{Say it:} \emph{tolvũ}
  \item \emph{Say it:} \emph{tolvũ}, once more
  \item \emph{Say it:} \emph{dukāndār cokhā tole che} — The shopkeeper weighs the rice.
  \item \emph{Recall:} say the Gujarati for shopkeeper, then the Gujarati for groceries, then say \emph{tolvũ} again
\end{itemize}

\begin{tcolorbox}[breakable,colback=teal!4,colframe=teal!35!black,title={Before you move on}]
What does \gu{તોલવું} mean? (\textquotedblleft{}To weigh\textquotedblright{}.) Say it once more.
\end{tcolorbox}
\section[\texorpdfstring{\gu{કિલો} (kilo)}{કિલો (kilo)}]{\gu{કિલો} (kilo) — kilo, kilogram}

\label{lesson:GU-C183-kilo}



\begin{quote}
Before the new one: say the Gujarati for groceries, then the Gujarati for to weigh.
\end{quote}

\subsection*{You'll want to know: \gu{કિલો}}
\textbf{\gu{કિલો}} — \emph{kilo} — \textquotedblleft{}kilo, kilogram\textquotedblright{}.

\begin{quote}
\textbf{\gu{એક} \gu{કિલો} \gu{ચોખા} \gu{સો} \gu{રૂપિયા} \gu{છે}.} — \emph{ek kilo cokhā so rūpiyā che} — \textquotedblleft{}A kilo of rice costs a hundred rupees.\textquotedblright{}
\end{quote}

A masculine noun.

\begin{cousinweb}
\textbf{\gu{કિલો}} (\emph{kilo}) is a loan from English \emph{kilo}, short for \emph{kilogram}.
\end{cousinweb}

\begin{grammarlens}[title={What you've built}]
One more word for everyday life at A2.
\end{grammarlens}

\subsection*{Your turn}
\begin{itemize}
\raggedright
  \item \emph{Say it:} \emph{kilo}
  \item \emph{Say it:} \emph{kilo}, once more
  \item \emph{Say it:} \emph{ek kilo cokhā so rūpiyā che} — A kilo of rice costs a hundred rupees.
  \item \emph{Recall:} say the Gujarati for groceries, then the Gujarati for to weigh, then say \emph{kilo} again
\end{itemize}

\begin{tcolorbox}[breakable,colback=teal!4,colframe=teal!35!black,title={Before you move on}]
What does \gu{કિલો} mean? (\textquotedblleft{}Kilo, kilogram\textquotedblright{}.) Say it once more.
\end{tcolorbox}
